% The figure atlas: every figure, in order, with the claim it makes.
%
% This is not the whitepaper. It exists so a change to `style.tex` can be
% reviewed in one place instead of four PDFs, and so a reader who will
% never build LaTeX can still see what the figures say.
%
% Build it after the figures: `make`, from the repository root.
\documentclass[10pt,a4paper]{article}

\usepackage[T1]{fontenc}
\usepackage[utf8]{inputenc}
\usepackage{lmodern}
\usepackage[a4paper,margin=2cm,landscape]{geometry}
\usepackage{graphicx}
\usepackage{xcolor}
\usepackage[hidelinks]{hyperref}

\graphicspath{{../../build/}}

\setlength{\parindent}{0pt}
\setlength{\parskip}{6pt}

\newcommand{\figurepage}[2]{%
  \clearpage
  \section*{#1}
  \vspace{-4pt}%
  #2\par
  \vspace{4pt}%
}

% Wide figures, one page each: scaled to the line, but never taller than
% the page, or LaTeX breaks the page in the middle of the picture.
\newcommand{\figimg}[1]{%
  \includegraphics[width=\linewidth,height=0.78\textheight,keepaspectratio]{#1}%
}

\begin{document}

\begin{center}
  {\Large\bfseries Workspace isolation}\\[2pt]
  {\large figures for the whitepaper}\\[6pt]
  {\small decouple-bench — draft, figures only}
\end{center}

\vspace{8pt}
These four figures answer one question: what does one run actually touch,
and what is kept apart from what. They describe the harness as it is, not
as it was meant to be — each file names the sources it was drawn from.

\figurepage{1. Two phases}{%
  \figimg{01-two-phases}%
}

\figurepage{2. The agent's boundary}{%
  \figimg{02-agent-boundary}%
}

\figurepage{3. Workspace provenance}{%
  \figimg{03-workspace-provenance}%
}

\figurepage{4. One case, one stack}{%
  \figimg{04-case-isolation}%
}

\end{document}
